\section{Цель работы и теоретические сведения}

Цель работы --- реализовать многоветвевое дерево решений для бинарной
классификации по категориальным признакам. Требуется оценить качество
на обучающей и тестовой выборках, построить ROC- и PR-кривые, вычислить
площади под ними методом трапеций и исследовать влияние порога на ошибки.
В дополнительном эксперименте сравниваются критерии энтропии и Джини.

Пусть \(S\) --- множество объектов в узле, а \(p_c\) --- доля класса \(c\).
Для каждого значения \(v\) признака \(A\) формируется подмножество \(S_v\).
Энтропия узла и информационный выигрыш разбиения равны
\[
H(S)=-\sum_{c=0}^{1}p_c\log_2 p_c,
\qquad
IG(S,A)=H(S)-\sum_v\frac{|S_v|}{|S|}H(S_v).
\]
Слагаемое \(p_c\log_2 p_c\) при \(p_c=0\) принимается равным нулю.
Для критерия Джини используется та же разность между неопределённостью
родительского узла и взвешенной неопределённостью потомков, где
\(G(S)=1-\sum_{c=0}^{1}p_c^2\) заменяет \(H(S)\).

В каждом узле выбирается доступный признак с наибольшим положительным
выигрышем. Выбранный признак исключается из дальнейших разбиений на
этом пути. Узел становится листом, если все объекты принадлежат одному
классу, признаки закончились или положительного выигрыша нет.
Оценка \(p_1\) берётся из достигнутого листа. При отсутствии ветви для
значения признака используется \(p_1\) текущего узла. Прогноз равен
единице при \(p_1\geq t\), где стандартный порог \(t=0{,}5\).

\section{Данные и разбиение}

Набор Higher Education Students Performance Evaluation содержит 145
наблюдений. Класс 1 соответствует \texttt{OUTPUT Grade} \(\geq 3\),
класс 0 --- остальным оценкам. \texttt{Student ID} и \texttt{OUTPUT Grade}
исключены из входных признаков. Все 31 оставшийся признак рассматриваются
как категориальные. Для пустых значений и «?» предусмотрена отдельная
категория «пропуск». В обработанных данных пропусков не оказалось.

\begin{table}[!ht]
\centering\small
\caption{Распределение объектов по выборкам}
\begin{tabular}{lrrrr}
\toprule
Выборка & Всего & Класс 0 & Класс 1 & Доля класса 1 \\
\midrule
Полная    & 145 & 67 & 78 & \(53{,}8\,\%\) \\
Обучающая & 116 & 54 & 62 & \(53{,}4\,\%\) \\
Тестовая  & 29  & 13 & 16 & \(55{,}2\,\%\) \\
\bottomrule
\end{tabular}
\end{table}

Объекты перемешаны внутри каждого класса с начальным состоянием
генератора 42. При разбиении 80/20 выделено 29 тестовых объектов, из
них 16 положительных по округлённой доле класса 1. Из 31 признака
один раз отобраны \(\lfloor\sqrt{31}\rfloor=5\) без возвращения
с начальным состоянием 43. В общий для обоих деревьев набор вошли
\texttt{Sex} (пол), \texttt{Additional work} (дополнительная работа),
\texttt{Accomodation type in Cyprus} (тип жилья),
\texttt{Preparation to midterm exams 1} (подготовка к промежуточным
экзаменам) и \texttt{Taking notes in classes} (конспектирование).

\clearpage
\section{Реализация дерева и метрик}

\subsection{Обучение и прогнозирование}

Функция \texttt{fit} строит дерево рекурсивно. Выигрыши в пределах
допуска \(10^{-12}\) считаются равными, поэтому при равенстве сохраняется
первый признак в порядке столбцов. Функция \texttt{route} возвращает узел,
чья доля класса 1 используется для прогноза.

\VerbatimInput[lastline=69,fontsize=\codefont,baselinestretch=0.96,numbers=left,numbersep=7pt,xleftmargin=19pt]{../src/tree.py}

\clearpage
\subsection{Метрики и площадь под кривой}

Для порога \(t\) вычисляются числа \(TP\), \(FP\), \(TN\) и \(FN\).
Точность положительных прогнозов обозначена \(P\), полнота --- \(R\).
\[
\mathrm{Accuracy}=\frac{TP+TN}{N},\qquad
P=\frac{TP}{TP+FP},\qquad
R=\frac{TP}{TP+FN},\qquad
\mathrm{FPR}=\frac{FP}{FP+TN}.
\]
При пустом наборе положительных прогнозов precision в таблице метрик
принимается равной нулю. Ордината начальной точки PR-кривой по
соглашению равна 1.

\VerbatimInput[fontsize=\codefont,numbers=left,numbersep=7pt,xleftmargin=19pt]{../src/metrics.py}

\clearpage
\section{Структура дерева и важность признаков}

Глубина энтропийного дерева равна 5 при нулевой глубине корня.
Дерево содержит 32 листа. Слева от стрелки указано значение входящей
ветви. Поле \texttt{n} содержит число обучающих объектов, \texttt{p1} ---
долю класса 1, \texttt{IG} --- информационный выигрыш. Ниже показаны узлы
до глубины 3, многоточием отмечены продолжения ветвей.

\begin{Verbatim}[fontsize=\small]
Корень -> Sex [IG=0.04546, n=116, p1=0.5345]
  1 -> Taking notes in classes [IG=0.06670, n=45, p1=0.3778]
    1 -> лист [n=3, p1=0.0000]
    2 -> Additional work [IG=0.08170, n=18, p1=0.5000]
      1 -> Accomodation type in Cyprus [IG=0.07278, n=9, p1=0.6667] ...
      2 -> Accomodation type in Cyprus [IG=0.16775, n=9, p1=0.3333] ...
    3 -> Additional work [IG=0.13955, n=24, p1=0.3333]
      1 -> Accomodation type in Cyprus [IG=0.19351, n=10, p1=0.1000] ...
      2 -> Preparation to midterm exams 1 [IG=0.25698, n=14, p1=0.5000] ...
  2 -> Preparation to midterm exams 1 [IG=0.05207, n=71, p1=0.6338]
    1 -> Additional work [IG=0.03978, n=56, p1=0.6429]
      1 -> Accomodation type in Cyprus [IG=0.04076, n=17, p1=0.4706] ...
      2 -> Taking notes in classes [IG=0.07504, n=39, p1=0.7179] ...
    2 -> Accomodation type in Cyprus [IG=0.29990, n=11, p1=0.4545]
      1 -> Additional work [IG=0.09236, n=8, p1=0.3750] ...
      2 -> лист [n=2, p1=1.0000]
      4 -> лист [n=1, p1=0.0000]
    3 -> лист [n=4, p1=1.0000]
\end{Verbatim}

\experimentfigure{0.97}{importance.png}{Суммарный информационный выигрыш по признакам}

Важность признака равна сумме его выигрышей во всех узлах, где он
использован, без дополнительного взвешивания по размеру узла.
Наибольшие значения дали тип жилья (\(1{,}129\)) и подготовка
к экзаменам (\(0{,}561\)). Признак \texttt{Sex} дал лучший выигрыш
в корне, однако использован только в одном узле.

\clearpage
\section{Качество классификации}

\begin{table}[!ht]
\centering\small
\caption{Метрики двух деревьев при пороге \(t=0{,}5\)}
\begin{tabular}{llrrrrrrr}
\toprule
Критерий & Выборка & \(TP\) & \(FP\) & \(TN\) & \(FN\) & Accuracy & Precision & Recall \\
\midrule
Энтропия & Обучение & 52 & 18 & 36 & 10 & 0,759 & 0,743 & 0,839 \\
Энтропия & Тест     &  8 &  3 & 10 &  8 & 0,621 & 0,727 & 0,500 \\
Джини    & Обучение & 52 & 18 & 36 & 10 & 0,759 & 0,743 & 0,839 \\
Джини    & Тест     &  8 &  3 & 10 &  8 & 0,621 & 0,727 & 0,500 \\
\bottomrule
\end{tabular}
\end{table}

При пороге \(0{,}5\) энтропийное дерево пропустило 8 из 16 объектов
класса 1 (ошибки второго рода). Ложноположительных ответов было 3
на 13 объектов класса 0 (ошибки первого рода).

\subsection{Кривые ROC и PR}

ROC строится в координатах \((\mathrm{FPR},R)\), PR ---
в координатах \((R,P)\). Порог последовательно принимает уникальные
доли класса 1 в листьях, значения 0 и 1, а также доли внутренних узлов,
фактически возвращённые на тесте при отсутствии ветви. Дополнительный
порог \(+\infty\) задаёт точки \((0,0)\) для ROC и \((0,1)\) для PR.
При пороге 0 к классу 1 относятся все объекты и ROC достигает
точки \((1,1)\). Горизонтальный пунктир на PR соответствует доле
положительного класса \(16/29\).

\experimentfigure{0.97}{curves.png}{ROC- и PR-кривые на тестовой выборке}

Площадь под каждой кривой вычислена суммированием площадей трапеций.
\[
\mathrm{AUC}=\sum_i (x_{i+1}-x_i)\frac{y_{i+1}+y_i}{2}.
\]
Перед интегрированием точки устойчиво сортируются по \(x\). При равных
абсциссах сохраняется порядок убывания порогов. Для энтропийного дерева
ROC-AUC равна \(0{,}668\), PR-AUC равна \(0{,}745\).

\clearpage
\subsection{Влияние порога и переобучение}

\begin{table}[!ht]
\centering\small
\caption{Ошибки энтропийного дерева при изменении порога на тесте}
\begin{tabular}{lrrrr}
\toprule
Порог & \(FP\) & \(FN\) & Precision & Recall \\
\midrule
\(1/3\) & 6 &  4 & 0,667 & 0,750 \\
0,50    & 3 &  8 & 0,727 & 0,500 \\
0,75    & 1 &  9 & 0,875 & 0,438 \\
1,00    & 1 & 12 & 0,800 & 0,250 \\
\bottomrule
\end{tabular}
\end{table}

При повышении порога с \(0{,}50\) до \(0{,}75\) число \(FP\) снизилось
с 3 до 1, а \(FN\) выросло с 8 до 9. При дальнейшем повышении порога
до \(1{,}00\) число \(FP\) осталось равным 1, а \(TP\) упало с 7 до 4,
поэтому precision снизилась с \(0{,}875\) до \(0{,}800\).
Снижение порога до \(1/3\) позволило найти 12 из 16 положительных
объектов при 6 ложных срабатываниях.

Accuracy энтропийного дерева снизилась с \(0{,}759\) на обучении до
\(0{,}621\) на тесте, то есть на 13,8 процентного пункта. Recall упала
с \(0{,}839\) до \(0{,}500\), а precision --- с \(0{,}743\) до \(0{,}727\).
Разрыв и листья, построенные по одному обучающему объекту, указывают
на переобучение. При 29 тестовых объектах величина разрыва
чувствительна к разбиению.

\section{Сравнение с критерием Джини}

Дерево Джини использует те же пять признаков, имеет глубину 5 и
32 листа. В обоих деревьях корень делит объекты по \texttt{Sex}.
Структура ветви \texttt{Sex} \(=1\) совпадает. При \texttt{Sex} \(=2\)
первое разбиение по Джини использует тип жилья, а по энтропии ---
подготовку к экзаменам. Для одного тестового объекта дерево Джини
вернуло долю класса 1 во внутреннем узле из-за отсутствующей ветви.

При пороге \(0{,}5\) метрики и числа ошибок совпали на обеих выборках.
У дерева Джини ROC-AUC равна \(0{,}666\), PR-AUC равна \(0{,}732\).
На этом разбиении энтропия дала площади больше примерно на \(0{,}002\)
и \(0{,}013\) соответственно.

\section{Заключение}

Наибольший вклад в разбиения энтропийного дерева внесли тип жилья и
подготовка к экзаменам. При стандартном пороге дерево обнаружило
половину положительных объектов тестовой выборки. Снижение порога
сокращает число пропусков ценой дополнительных ложных срабатываний.
Падение accuracy и recall на тесте указывает на ограниченную
переносимость обучающих разбиений. Замена энтропии на критерий Джини
изменила часть ветвей, сохранив разрыв между обучением и тестом.
